\section{Récupération d’un modèle de RAM}

\subsection{Validation du fonctionnement de la RAM}

Nous récupérons le modèle \texttt{ramchip.vhd} et son testbench afin de pouvoir observer et analyser son comportement sur ModelSim.

% Injecter le graphe de fonctionnement de la RAM
\begin{figure}[H]
    \centering
    \includegraphics[width=1\textwidth]{1_RAM_modified.png}
    \caption{Chronogramme de la RAM}
    \label{fig:ram}
\end{figure}

Dans le chronogramme ci-dessus, nous observons que notre simulation représente correctement le comportement d'une RAM\@.
\\

Le bon fonctionnement se valide en activant \texttt{we}. À chaque front montant de \texttt{clock}, nous définissons l’adresse \texttt{address} de la donnée à écrire, puis nous plaçons cette donnée sur \texttt{datain}.
\\

\begin{itemize}
    \item À 20 ns : écriture de 01 à l’adresse 001 ; \texttt{dataout} restitue 01
    \item À 30 ns : écriture de 02 à l’adresse 002 ; \texttt{dataout} restitue 02
    \item À 40 ns : écriture de 03 à l’adresse 003 ; \texttt{dataout} restitue 03, et ainsi de suite
    \\
\end{itemize}

Un effet de décalage apparait, la valeur affichée correspond toujours à la donnée appliquée au cycle précédent.
Exemple entre $20\text{ ns}$ et $30\text{ ns}$, \texttt{address} vaut 001, \texttt{datain} vaut 01 et \texttt{dataout} restitue 00. Cet effet est normal car notre RAM est synchrone et ne restitue donc la donnée qu'au cycle suivant.
\\


\subsection{Modification du code de la RAM}

Dans un second temps, nous mettons à jour le code en remplaçant les \texttt{std\_logic\_vector} par des \texttt{unsigned} 
afin de pouvoir effectuer des opérations arithmétiques sur les adresses. Le fonctionnement reste identique à celui de la RAM initiale.
\\

\subsection{Ajout de la généricité de la RAM}

Pour rendre notre RAM plus générique, nous ajoutons des paramètres de taille de données et d'adresses.
Nous ajoutons donc les généricités \texttt{DATA\_LENGHT} et \texttt{ADDR\_LENGHT} qui permettent de définir la taille de la RAM\@.

\begin{lstlisting}[caption={Nouvelle entité de la RAM}, label={lst:ram_entity}]
entity sync_ram is
  generic (
    ADDR_LENGHT : integer := 10;
    DATA_LENGHT : integer := 8
  );
  port (
    clock   : in  std_logic;
    we      : in  std_logic;
    address : in  unsigned((ADDR_LENGHT - 1) downto 0);
    datain  : in  unsigned((DATA_LENGHT - 1) downto 0);
    dataout : out unsigned((DATA_LENGHT - 1) downto 0)
  );
end entity sync_ram;
\end{lstlisting}

\subsection{Ajout du signal de sélection}

L'ajout du signal \texttt{cs} permet de contrôler l'accès à la RAM et de placer sa sortie en haute impédance (\texttt{'Z'}) lorsqu'elle n'est pas sélectionnée, en vue de la création de la \texttt{quadram}.

\begin{lstlisting}[caption={Ajout du signal cs}, label={lst:ram_entity_cs}]
entity sync_ram is
  generic(DataLenght : INTEGER := 8; AddrLenght : INTEGER := 8);
  port (
    cs    : in std_logic;
    clock : in std_logic;
    ...
  );
end entity;

architecture RTL of sync_ram is
  ...
begin
  process(clock) is
  begin
    if rising_edge(clock) and cs = '1' then
      -- Ecriture synchronisee
    end if;
  end process;

  dataout <= ram(to_integer(unsigned(read_address))) when cs = '1' else (others => 'Z');
end architecture;
\end{lstlisting}

\begin{figure}[H]
    \centering
    \includegraphics[width=1\textwidth]{1_RAM_CS_modified.png}
    \caption{Chronogramme de la RAM avec le signal \texttt{cs}}
    \label{fig:ram_cs}
\end{figure}

La Figure~\ref{fig:ram_cs} présente le chronogramme de la RAM après l’ajout du signal de sélection \texttt{cs}. Lorsque \texttt{cs} vaut \texttt{'1'}, la RAM est sélectionnée : l’écriture est autorisée sur un front montant de \texttt{clock} si \texttt{we} est également actif, et la sortie \texttt{dataout} restitue la donnée mémorisée à l’adresse demandée. Lorsque \texttt{cs} vaut \texttt{'0'}, la mémoire n’est pas sélectionnée et sa sortie est placée en haute impédance (\texttt{'Z'}), afin de ne pas perturber un éventuel bus partagé. \\

Le chronogramme illustre le comportement synchrone de l’écriture : les données de \texttt{datain} sont enregistrées sur un front montant de \texttt{clock}, lorsque \texttt{cs} et \texttt{we} sont actifs. La sortie \texttt{dataout} est, quant à elle, conditionnée par \texttt{cs}. L’ajout de ce signal complète donc le fonctionnement de la RAM en permettant de contrôler son accès.